\documentclass[12pt]{article}
\usepackage[utf8]{inputenc}
\usepackage[T1]{fontenc}
\usepackage{hyperref}
\title{Toda Language Resources for K-12 Curriculum Development}
\author{Compiled by Claude Sonnet 5 (claude-sonnet-5) for Andr\'as Kornai}
\date{}
\begin{document}
\maketitle

\section{Introduction}
Toda (ISO 639-3 tcx; Glottocode toda1252) is a Dravidian language of the Nilgiri hills in India, with about 1,600 speakers. The reference description says it is stable and not known to be taught in schools. Search found scholarly grammars but no school materials and no dedicated Toda dictionary.

\section{Primary Educational Resources}
None found.

\section{Linguistic and Reference Materials}
\begin{itemize}
\item S. Sakthivel, \emph{A Grammar of the Toda Language}, Annamalai University, Centre of Advanced Study in Linguistics, Publication No.~49, 1977 (PDF): \url{https://www.tamildigitallibrary.in/admin/assets/book/TVA_BOK_0018991_Toda_Languag.pdf}
\item M.\,B. Emeneau, \emph{Toda Grammar and Texts} (1984) and \emph{Toda Songs} (1971). These are cited as standard references in the French study at: \url{https://books.openedition.org/ifp/7716}
\item A bibliography for the study of Nilgiri hill tribes, which lists Emeneau's unpublished fieldnotes on Toda (Berkeley, 1935--38): \url{https://www.tamildigitallibrary.in/admin/assets/book/TVA_BOK_0019063_A_Bibliography_for_the_study_of_nilagiri_hill_tribes.pdf}
\item Wikipedia overview of Toda (a starting point only): \url{https://wikipedia.com/wiki/ISO_639:tcx}
\end{itemize}

\section{Religious and Community Resources}
None found.

\section{Digital and Audio Resources}
None found in search.

\section{Curriculum Development Context}
Toda has strong community use but little published pedagogical material. The grammars and Emeneau's texts could serve as a base for orthography and materials development with the community.
\end{document}